% Options for packages loaded elsewhere
\PassOptionsToPackage{unicode}{hyperref}
\PassOptionsToPackage{hyphens}{url}
\documentclass[
  a4paper,
]{ltjsarticle}
\usepackage{xcolor}
\usepackage[margin=25mm]{geometry}
\usepackage{amsmath,amssymb}
\setcounter{secnumdepth}{-\maxdimen} % remove section numbering
\usepackage{iftex}
\ifPDFTeX
  \usepackage[T1]{fontenc}
  \usepackage[utf8]{inputenc}
  \usepackage{textcomp} % provide euro and other symbols
\else % if luatex or xetex
  \usepackage{unicode-math} % this also loads fontspec
  \defaultfontfeatures{Scale=MatchLowercase}
  \defaultfontfeatures[\rmfamily]{Ligatures=TeX,Scale=1}
\fi
\usepackage{lmodern}
\ifPDFTeX\else
  % xetex/luatex font selection
\fi
% Use upquote if available, for straight quotes in verbatim environments
\IfFileExists{upquote.sty}{\usepackage{upquote}}{}
\IfFileExists{microtype.sty}{% use microtype if available
  \usepackage[]{microtype}
  \UseMicrotypeSet[protrusion]{basicmath} % disable protrusion for tt fonts
}{}
\makeatletter
\@ifundefined{KOMAClassName}{% if non-KOMA class
  \IfFileExists{parskip.sty}{%
    \usepackage{parskip}
  }{% else
    \setlength{\parindent}{0pt}
    \setlength{\parskip}{6pt plus 2pt minus 1pt}}
}{% if KOMA class
  \KOMAoptions{parskip=half}}
\makeatother
\usepackage{longtable,booktabs,array}
\newcounter{none} % for unnumbered tables
\usepackage{calc} % for calculating minipage widths
% Correct order of tables after \paragraph or \subparagraph
\usepackage{etoolbox}
\makeatletter
\patchcmd\longtable{\par}{\if@noskipsec\mbox{}\fi\par}{}{}
\makeatother
% Allow footnotes in longtable head/foot
\IfFileExists{footnotehyper.sty}{\usepackage{footnotehyper}}{\usepackage{footnote}}
\makesavenoteenv{longtable}
\usepackage{graphicx}
\makeatletter
\newsavebox\pandoc@box
\newcommand*\pandocbounded[1]{% scales image to fit in text height/width
  \sbox\pandoc@box{#1}%
  \Gscale@div\@tempa{\textheight}{\dimexpr\ht\pandoc@box+\dp\pandoc@box\relax}%
  \Gscale@div\@tempb{\linewidth}{\wd\pandoc@box}%
  \ifdim\@tempb\p@<\@tempa\p@\let\@tempa\@tempb\fi% select the smaller of both
  \ifdim\@tempa\p@<\p@\scalebox{\@tempa}{\usebox\pandoc@box}%
  \else\usebox{\pandoc@box}%
  \fi%
}
% Set default figure placement to htbp
\def\fps@figure{htbp}
\makeatother
\setlength{\emergencystretch}{3em} % prevent overfull lines
\providecommand{\tightlist}{%
  \setlength{\itemsep}{0pt}\setlength{\parskip}{0pt}}
% Glyph map for lualatex (newunicodechar). Maps symbols that may lack font glyphs.
\usepackage{newunicodechar}
\usepackage{amssymb}
\newunicodechar{∎}{\ensuremath{\blacksquare}}
\newunicodechar{✓}{\ensuremath{\checkmark}}
\newunicodechar{✗}{\ensuremath{\times}}
\newunicodechar{⟹}{\ensuremath{\Longrightarrow}}
\newunicodechar{⟺}{\ensuremath{\Longleftrightarrow}}
\newunicodechar{↪}{\ensuremath{\hookrightarrow}}
\newunicodechar{⌊}{\ensuremath{\lfloor}}
\newunicodechar{⌋}{\ensuremath{\rfloor}}
\newunicodechar{≃}{\ensuremath{\simeq}}
\newunicodechar{≈}{\ensuremath{\approx}}
\newunicodechar{≤}{\ensuremath{\leq}}
\newunicodechar{≥}{\ensuremath{\geq}}
\newunicodechar{≠}{\ensuremath{\neq}}
\newunicodechar{→}{\ensuremath{\rightarrow}}
\newunicodechar{←}{\ensuremath{\leftarrow}}
\newunicodechar{↔}{\ensuremath{\leftrightarrow}}
\newunicodechar{⊥}{\ensuremath{\perp}}
\newunicodechar{√}{\ensuremath{\surd}}
\newunicodechar{∑}{\ensuremath{\sum}}
\newunicodechar{∏}{\ensuremath{\prod}}
\newunicodechar{∞}{\ensuremath{\infty}}
\newunicodechar{∈}{\ensuremath{\in}}
\newunicodechar{∀}{\ensuremath{\forall}}
\newunicodechar{∣}{\ensuremath{\mid}}
\newunicodechar{γ}{\ensuremath{\gamma}}
\newunicodechar{λ}{\ensuremath{\lambda}}
\newunicodechar{μ}{\ensuremath{\mu}}
\newunicodechar{ρ}{\ensuremath{\rho}}
\newunicodechar{σ}{\ensuremath{\sigma}}
\newunicodechar{φ}{\ensuremath{\varphi}}
\newunicodechar{ψ}{\ensuremath{\psi}}
\newunicodechar{θ}{\ensuremath{\theta}}
\newunicodechar{Δ}{\ensuremath{\Delta}}
\newunicodechar{Π}{\ensuremath{\Pi}}
\newunicodechar{Λ}{\ensuremath{\Lambda}}
\newunicodechar{τ}{\ensuremath{\tau}}
\newunicodechar{ε}{\ensuremath{\varepsilon}}
\newunicodechar{ω}{\ensuremath{\omega}}
\newunicodechar{π}{\ensuremath{\pi}}
\newunicodechar{ν}{\ensuremath{\nu}}
\newunicodechar{±}{\ensuremath{\pm}}
\newunicodechar{×}{\ensuremath{\times}}
\newunicodechar{·}{\ensuremath{\cdot}}
\usepackage{bookmark}
\IfFileExists{xurl.sty}{\usepackage{xurl}}{} % add URL line breaks if available
\urlstyle{same}
\hypersetup{
  pdftitle={The Qualification for Ignition Is an Unstable Relative Equilibrium --- A Qualification Screening of Three Families of Initial Values for Relational-Wave Closed Systems},
  pdfauthor={Noriaki Kihara (WF System Co., Ltd.), ORCID 0009-0004-6753-4020},
  hidelinks,
  pdfcreator={LaTeX via pandoc}}

\title{The Qualification for Ignition Is an Unstable Relative
Equilibrium --- A Qualification Screening of Three Families of Initial
Values for Relational-Wave Closed Systems}
\author{Noriaki Kihara (WF System Co., Ltd.), ORCID 0009-0004-6753-4020}
\date{2026-09-12}

\begin{document}
\maketitle

\textbf{Series:} Foundations of Self-Consistent Relational-Wave Closed
Systems (Paper 2 of 10)\\
\textbf{Author:} Noriaki Kihara (WF System Co., Ltd.)　\textbf{Date:}
2026-09-12\\
\textbf{Version DOI:} 10.5281/zenodo.22728824\\
\textbf{Concept DOI:} 10.5281/zenodo.22728823\\
\textbf{ORCID:} 0009-0004-6753-4020\\
\textbf{Zenodo:} https://zenodo.org/records/22728824\\

\begin{center}\rule{0.5\linewidth}{0.5pt}\end{center}

\subsection{Abstract}\label{abstract}

The inflation (Paper 6) and SSB (Paper 5) of relational-wave closed
systems start from a specific initial value (floor). We establish,
through a controlled experiment that takes the initial value as the
control variable, that this initial value is not a ``conveniently
planted seed.'' Leaving the physical map unchanged, we substituted only
the initial value with three families------(A) symmetric parents with
phases \(0^\circ/90^\circ\) (equal-amplitude system), (B) zero-centroid
parents, (C) high-symmetry theoretical floor (cyclic-Fourier relative
equilibrium \(z_{ij}=r_d e^{i2\pi(i+j)/N}\))------and ran 500 steps for
\(N=3,\ldots,40\). Results: (A) is a neutral fixed point for
\(N\equiv1\ (\mathrm{mod}\,4)\), and for other \(N\) it relaxes
immediately at step 1 (no ignition); (B) also relaxes immediately with
residual \(O(1)\); \textbf{only (C) has the step-1 jump vanish in all 38
systems, igniting via an exponential ramp from the floor} (onset
\(65\)--\(421\) for \(N=4\)--\(17\); for \(N\ge18\) a slow ramp is in
progress; \(N=3\) is a neutral floor). \textbf{The qualification for
ignition is ``an unstable relative equilibrium (a self-consistent
eigenmode \(H(z)z=\lambda z\), residual at machine precision)''; high
symmetry, zero centroid, and square-closure are not qualifications.} The
failures of (A)(B) are established as arising not from ``being highly
symmetric'' but from ``not having been a relative equilibrium.'' The
floor data used by Papers 5, 6, and 7 is the make\_parent eigenmode
floor (Experiment 13) that satisfies this qualification. Classification:
hypothesis → refutation → established (controlled experiment).

\begin{center}\rule{0.5\linewidth}{0.5pt}\end{center}

\subsection{1. The Problem------A Response to ``the Conveniently Planted
Seed''}\label{the-problema-response-to-the-conveniently-planted-seed}

If the floor were conveniently chosen for the result, the very rapid
expansion of this system would be suspect. Therefore we need to
disentangle, through a controlled experiment that takes the initial
value as the sole control variable, \textbf{which of the properties of
the initial value determines ignition}. The candidate ``seemingly good
properties'' are high symmetry, zero centroid, square-closure
\(\sum z^2=0\), and relative equilibrium (fixed point). This paper
separates these to identify the qualification.

\subsection{2. The Dynamics and the Qualification of Relative
Equilibrium}\label{the-dynamics-and-the-qualification-of-relative-equilibrium}

The map is the one\_step \(z\to\exp(\Delta\tau K)z\),
\(K_{ef}=A_{ef}\sin(\varphi_f-\varphi_e)\), \(\Delta\tau=2\pi/N\) of the
canonical source \texttt{run\_N3\_N40\_stage123\_v1.py} (SHA256
\texttt{1abf2353…}) (Paper 1 §2). For the floor to be a relative
equilibrium that ``stays at rest when placed'' means that \(z\) is a
self-consistent eigenmode \(H(z)z=\lambda z\) (\(H=iK\)) of the
generator. If it is a relative equilibrium, no immediate jump occurs at
step 1, and if the equilibrium is \textbf{unstable}, the infinitesimal
component is amplified exponentially and ignites (Paper 6). \textbf{We
define the qualification for ignition as follows: the floor is a
relative equilibrium (a self-consistent eigenmode) and possesses a
linearly unstable direction (the floor Jacobian has
\(\lambda_{\max}>0\), or a positive exponential growth rate is
measured).} This is a \textbf{separate concept} from ``whether the onset
threshold was crossed within the observation time (500 steps in this
experiment)''; even with the qualification, if the clock is slow, it can
remain uncrossed within the observation window (the \(N\ge18\) of §4).
The qualification is established up to the point of departure from the
floor (exponential amplification of the unstable direction); whether it
locks at the terminus is a separate dynamical fact (Paper 5). The run
changes only PARENT\_DIR/OUT of the physical canonical source (with SHA
match confirmed, the physical equations unchanged); the decision rules
are not put into the dynamics but only into post-run analysis.

\subsection{3. The Three Families of Initial Values Plus the Data
Floor}\label{the-three-families-of-initial-values-plus-the-data-floor}

{\def\LTcaptype{none} % do not increment counter
\begin{longtable}[]{@{}
  >{\raggedright\arraybackslash}p{(\linewidth - 10\tabcolsep) * \real{0.1667}}
  >{\raggedright\arraybackslash}p{(\linewidth - 10\tabcolsep) * \real{0.1667}}
  >{\raggedright\arraybackslash}p{(\linewidth - 10\tabcolsep) * \real{0.1667}}
  >{\raggedright\arraybackslash}p{(\linewidth - 10\tabcolsep) * \real{0.1667}}
  >{\raggedright\arraybackslash}p{(\linewidth - 10\tabcolsep) * \real{0.1667}}
  >{\raggedright\arraybackslash}p{(\linewidth - 10\tabcolsep) * \real{0.1667}}@{}}
\toprule\noalign{}
\begin{minipage}[b]{\linewidth}\raggedright
Family
\end{minipage} & \begin{minipage}[b]{\linewidth}\raggedright
Phase
\end{minipage} & \begin{minipage}[b]{\linewidth}\raggedright
Amplitude
\end{minipage} & \begin{minipage}[b]{\linewidth}\raggedright
Centroid \(\sum z\)
\end{minipage} & \begin{minipage}[b]{\linewidth}\raggedright
Square-closure
\end{minipage} & \begin{minipage}[b]{\linewidth}\raggedright
Relative equilibrium?
\end{minipage} \\
\midrule\noalign{}
\endhead
\bottomrule\noalign{}
\endlastfoot
(A) Symmetric parent v2 & \(0^\circ/90^\circ\) only & Family-unified
(even \(M\) equal-amplitude / odd \(M\) 2-valued) & Largest-class
\(\sqrt{M/2}\) & \(\approx0\) & Only \(N\equiv1\,(4)\) (neutral) \\
(B) Zero-centroid v3 & 4-position ± pairs/trios & --- & Exactly zero &
\(\approx0\) & No (residual \(O(1)\)) \\
(C) High-symmetry theoretical floor & \(2\pi(i+j)/N\) (cyclic Fourier) &
Distance-class constant \(r_d\) & \(\le2.0\times10^{-15}\) &
\(\le3.1\times10^{-16}\) & \textbf{Yes} (residual
\(4.1\text{e-}16\)--\(4.8\text{e-}14\)) \\
Data floor make\_parent & \(0/90\) lattice (Z4 residual
\(\le2\text{e-}12\)) & Per-edge (self-consistent adjustment) & Nonzero
\(0.25\)--\(1.41\) & \(\approx0\) (\(P_A=P_B\)) & \textbf{Yes} (residual
\(\sim10^{-13}\)) \\
\end{longtable}
}

\begin{enumerate}
\def\labelenumi{(\Alph{enumi})}
\tightlist
\item
  becomes an exact relative equilibrium (\(\sigma=N-1\)) only for
  two-axis equal amplitude with \(n_A=n_B\) (the \((N-1)/2\)-regular
  graph of \(N\equiv1\ \mathrm{mod}\,4\)), because
  \(P_A=n_A r^2=P_B=n_B r^2\) is required for closure; otherwise it is
  not a fixed point. The generation of (C) follows a derivation
  procedure (fix phases → contract distance classes → eigenvalue problem
  → full-space verification) and does not use make\_parent.
\end{enumerate}

\subsection{4. Results of the Qualification Screening (Experiment 12,
Experiment
13)}\label{results-of-the-qualification-screening-experiment-12-experiment-13}

\begin{itemize}
\tightlist
\item
  \textbf{(A) Symmetric parent v2 is NG}: only \(N=3\) ramps
  (gauge-identical to the old make\_parent floor),
  \(N\equiv1\ \mathrm{mod}\,4\) is a neutral fixed point at the floor
  (no ignition), and for other \(N\) it \textbf{relaxes immediately} at
  step 1 to \(f(1)=7.5\times10^{-4}\)--\(2.7\times10^{-2}\) with no
  linear phase. Decomposition: \(\sum z^2\) held, but the cross-term
  channel was empty (all waves \(ab\equiv0\)) and the centroid was
  largest-class.
\item
  \textbf{(B) Zero-centroid v3 is also NG} (trial runs \(N=3\)--\(6\)):
  even satisfying the 4 conditions (phase, zero centroid,
  square-closure, normalization) exactly, because it is not a relative
  equilibrium (residual \(O(1)\)) it relaxes immediately at step 1 to
  \(0.15\)--\(0.52\).
\item
  \textbf{(C) The high-symmetry theoretical floor is OK}: the immediate
  jump at step 1 \textbf{vanishes in all 38 systems}
  (\(f(1)=5\times10^{-32}\)--\(8\times10^{-31}\)). \(N=4\)--\(17\)
  ignites via an exponential ramp from the floor (onset \(>0.05\) is
  \(65\) (\(N=4\))--\(421\) (\(N=17\))), and the amplification rate of
  \(N=4\), \(0.470\ \log_{10}/\text{step}\), nearly matches the old
  floor's \(0.475\) and the saturation \(1/6\) also matches. \(N\ge18\)
  has a positive amplification rate
  (\(0.060\)--\(0.015\ \log_{10}/\text{step}\), monotonically
  decreasing) with the 500 steps running out while the ramp is in
  progress = \textbf{unstable in all systems (qualified for ignition),
  only the clock is slow with increasing \(N\)}. \(N=3\) is a neutral
  floor (\(\lambda=-3/2\), consistent with the \(N=3\) final-state lock
  value). \(H_{\rm total}=1.000000\) (unitary conservation), global
  closure maximum \(\sim10^{-13}\).
\end{itemize}

\subsection{5. Established------The Qualification Is an ``Unstable
Relative
Equilibrium''}\label{establishedthe-qualification-is-an-unstable-relative-equilibrium}

The hypothesis (that one of zero centroid / square-closure / high
symmetry is the qualification for ignition) is refuted:

\begin{itemize}
\tightlist
\item
  \textbf{Zero centroid is not a qualification}: (B) does not ignite
  even with the centroid exactly zero. On the other hand, the data floor
  make\_parent ignites even with a nonzero centroid
  (\(0.25\)--\(1.41\)). The value of the centroid does not determine the
  presence or absence of ignition.
\item
  \textbf{Square-closure is not a qualification}: (A)(B) do not ignite
  even with \(\sum z^2\approx0\).
\item
  \textbf{High symmetry itself is not a qualification}: (A) is the most
  symmetric but does not ignite.
\end{itemize}

\textbf{Established: the qualification for ignition is solely ``an
unstable relative equilibrium (a self-consistent eigenmode, residual at
machine precision).''} The failures of (A)(B) are not because they are
``highly symmetric'' but because they ``were not relative equilibria.''
This independently re-confirms, through an initial-value controlled
experiment, the ``amplification ⟺ unstable relative equilibrium'' of the
onset-mode discrimination paper {[}2{]}. Therefore the floor used by
Papers 5 and 6 is not a ``seed that conveniently moves'' but an unstable
starting point determined solely by the map plus the selection criterion
``the most symmetric relative equilibrium'' (that a self-consistent
floor can be built on the 90-degree skeleton is Paper 1 (\(D^{-1}KD=iW\)
· Perron); the algebraic exact solution of the 90-degree skeleton is
Paper 3; the analytical backing for why the floor becomes 90-degree is
Paper 0).

From this controlled experiment, initial states are dynamically divided
into three onset modes. A state that is not a relative equilibrium
relaxes immediately, a stable or neutral relative equilibrium stays at
the floor, and only an unstable relative equilibrium departs
exponentially after a latency period. (B) is not a relative equilibrium
(immediate relaxation), (A) is a relative equilibrium but neutral (stays
at the floor), and (C) and make\_parent are unstable relative equilibria
(exponential ramp after latency)------the three families actually
separate this three-way classification. Therefore ``ignition'' is
classified not by the apparent symmetry or closure of the initial value
but by the linear stability of the relative equilibrium. \textbf{The
stability class of the initial state determines the onset mode of the
time evolution.}

\subsection{6. The Floor Used by the Paper Data (make\_parent Eigenmode
Floor)}\label{the-floor-used-by-the-paper-data-make_parent-eigenmode-floor}

The run data of Papers 5, 6, 7, and 8 (\texttt{full\_N3\_N40\_sweep})
uses a floor obtained by finding the self-consistent eigenmode \(v\)
through the \texttt{make\_parent} iteration of the canonical engine
(\texttt{rng=default\_rng(40260722+1000N)}, \texttt{iters=1200},
\texttt{tol=1e-12}), and setting \(Z_0=(v+10^{-15}g)/\|\cdot\|\) (\(g\)
= the infinitesimal seed of the zero-closure kernel). Structural audit
(all 38 systems): Z4 phase residual \(\le2.0\times10^{-12}\) rad,
\(P_{\rm even}=P_{\rm odd}=0.500000\) exactly, Rayleigh residual
\(\sim10^{-15}\)--\(10^{-13}\). The \(v,g,Z_0\) of \(N=40\) are
bit-identical to the July canonical static parent (control gate). The
step-0 complex plane (all \(N=3\)--\(40\)) is shown in Figure 1------it
is a 90-degree lattice in which every wave lies on one of two orthogonal
axes (a \textbf{different member} that algebraically and exactly
constructs a self-consistent floor on the 90-degree skeleton------with
the same ignition qualification but a different \(W\), amplitude, and
centroid------is treated in Paper 3. It is not an exact version of the
make\_parent floor itself, Paper 0).

\begin{figure}
\centering
\pandocbounded{\includegraphics[keepaspectratio,alt={Figure 1: make\_parent floor step-0 complex plane (all N=3..40). A 90-degree phase lattice in which every wave lies on two orthogonal axes.}]{figs/p2_ignition_qualification_en_fig1.png}}
\caption{Figure 1: make\_parent floor step-0 complex plane (all
N=3..40). A 90-degree phase lattice in which every wave lies on two
orthogonal axes.}
\end{figure}

\textbf{Figure 1} The step-0 complex plane of the data floor
(make\_parent) (\(N=3\)--\(40\)).

\subsection{7. Claim Classification and Open
Items}\label{claim-classification-and-open-items}

\begin{itemize}
\tightlist
\item
  Classification: hypothesis → refutation → established (a controlled
  experiment taking the initial value as the sole control variable).
  Decision: the qualification = unstable relative equilibrium (relative
  equilibrium + linearly unstable direction) is established.
\item
  The qualification and the firing within the observation window are
  separate (§2): with the (C) theoretical floor, \(N=4\)--\(17\) crossed
  the onset threshold within 500 steps; \(N\ge18\) has a positive
  amplification rate (\(0.060\)--\(0.015\,\log_{10}/\text{step}\)) and
  is \textbf{qualified for ignition = departs from the floor}, but the
  clock is slow and it is uncrossed within 500 steps; \(N=3\) is neutral
  (\(\lambda=-3/2\), unqualified). What this paper establishes is up to
  this qualification (departure from the floor); \textbf{whether or not
  it locks at the terminus is out of scope and is treated by Paper 5}
  (with the data floor, the lock of \(\Delta_{\rm tail}=-2\pi/N\) ·
  equal amplitude has been confirmed up to \(N=64\)).
\item
  Open: the moduli of the relative equilibrium (with a different
  generating measure a different self-consistent fixed point exists; the
  90-degree floor is the most symmetric member). The origin of the
  neutrality of \(N=3\) (\(\lambda=-3/2\)).
\end{itemize}

\subsection{Related work (structural agreement, not the source of
derivation)}\label{related-work-structural-agreement-not-the-source-of-derivation}

\begin{itemize}
\tightlist
\item
  Onset-mode discrimination paper (amplification ⟺ unstable relative
  equilibrium · all 86 states) Concept DOI 10.5281/zenodo.21798854: this
  paper is its initial-value-controlled version.
\item
  Relative equilibrium (Marsden \& Ratiu, geometric mechanics): the
  floor = a rigid-rotation relative equilibrium.
\end{itemize}

\subsection{References}\label{references}

\begin{enumerate}
\def\labelenumi{\arabic{enumi}.}
\tightlist
\item
  Noriaki Kihara, ``The mechanism of inflationary rapid expansion in
  self-consistent relational-wave closed systems,'' Concept DOI
  10.5281/zenodo.22112008.
\item
  Noriaki Kihara, ``Onset-mode discrimination (amplification and
  unstable relative equilibrium),'' Concept DOI 10.5281/zenodo.21798854.
\item
  J. E. Marsden and T. S. Ratiu, \emph{Introduction to Mechanics and
  Symmetry}, Springer (1999).
\end{enumerate}

\subsection{Reproduction}\label{reproduction}

Controlled experiment (three families):
\texttt{対称親v2\_500step走行\_20260906/}
(\texttt{make\_parents\_theoretical\_floor\_v1.py} = generation of (C),
\texttt{make\_parents\_v3\_centroid\_zero\_v1.py} = (B),
\texttt{wrapper\_run\_symmetric500\_v1.py},
\texttt{再実験結果メモ\_理論床\_N3\_N40.md},
\texttt{compare\_theoretical\_floor\_vs\_old\_makeparent.csv},
\texttt{theoretical\_floor\_parent\_manifest.json},
\texttt{SHA256SUMS.txt}). (A) symmetric parent is in
\texttt{最も対称性の高い初期値\_20260906/}. The data floor make\_parent
and its structural audit and bit reproduction:
\texttt{make\_parent型初期値\_自己無撞着構造\_20260907/}
(\texttt{make\_static\_parents\_N3\_N40\_v1.py},
\texttt{audit\_parent\_structure\_N3\_N40\_v1.py},
\texttt{full\_N3\_N40\_sweep/}). The run changes only PARENT\_DIR/OUT of
the physical canonical source \texttt{run\_N3\_N40\_stage123\_v1.py}
(SHA256
\texttt{1abf2353fee2e4f56f05e7a6f149fd086885136beb61ab571b48a56b09691567}),
with the physical equations unchanged. * The data npz files are saved in
each folder due to size (regenerable with SHA256SUMS + run\_all.sh).

\end{document}
